\documentclass[10pt,a4paper]{amsart}
\usepackage[margin=19mm]{geometry}
\usepackage{amsmath,amssymb,mathtools}
\usepackage[scr=rsfs,cal=euler]{mathalfa}
\usepackage{xfrac}
\usepackage[unicode,hidelinks,bookmarks=false]{hyperref}
\newcommand{\ie}{i.e.\ }
\newcommand{\cf}{cf.\ }
\newcommand{\resp}{respectively\ }
\newcommand{\Subsection}{\subsection}
\providecommand{\nobreakdash}{\nobreak\hbox{-}}
\setlength{\emergencystretch}{2em}
\multlinegap0pt
\usepackage{algorithm}\usepackage{algpseudocode}
\usepackage{amssymb}
\usepackage{float}
\newtheorem{theorem}{Theorem}
\newtheorem{assumption}{Assumption}
\newtheorem{lemma}{Lemma}
\title{T-SKM-Net: Trainable Neural Network Framework for Linear Constraint Satisfaction via Sampling Kaczmarz-Motzkin Method}
\author{Haoyu Zhu, Yao Zhang, Jiashen Ren, Qingchun Hou}
\date{Source: arXiv:2512.10461v1 (CC BY 4.0)}
\begin{document}
\maketitle

\noindent\emph{Source and licence.} T-SKM-Net: Trainable Neural Network Framework for Linear Constraint Satisfaction via Sampling Kaczmarz-Motzkin Method. Version arXiv:2512.10461v1 (CC BY 4.0), distributed by arXiv under the Creative Commons Attribution 4.0 International licence, http://creativecommons.org/licenses/by/4.0/, as recorded in the arXiv licence metadata for this exact version and shown on the abstract page https://arxiv.org/abs/2512.10461v1. Full licence notice: \texttt{licenses/arxiv-2512.10461-CC-BY-4.0.txt}.\par\smallskip

\noindent\textit{Editorial scope.} Counted unit: the article's theorem thm:appendix-trainability (source0.tex lines 868--878). The article's complete original proof of it is delivered with it (source0.tex lines 880--939, one \texttt{proof} environment of 60 lines). No mathematics is written, repaired or re-derived: the statement and the proof are the article's own lines, in the article's own notation, and no retained line is substituted or re-flowed.  Retained uncounted as the source-local context the proof needs: lines 720--725; lines 843--846; lines 863--866.  Nothing was omitted from any retained span, so the package carries no omission record.  No retained line contains a citation, so the wrapper carries no bibliography and no bibliography entry.  No retained line contains a figure environment, an image-inclusion command or drawing code, so no image is delivered and the package declares no asset dependency.  Editorial formatting only, in addition: an AMS \texttt{amsart} preamble that restates the article's own declarations and macros used by the retained lines, namely the environments \texttt{theorem}, \texttt{assumption}, \texttt{lemma} on separate counters, together with the standard packages those lines need.  The retained environments print in this document as Theorem 1, Assumption 1, Lemma 1 in order of appearance, whereas the article prints the same items with the numbering of its own full document.  The complete native source of the article as distributed in the arXiv e-print for this exact version is bundled unchanged as \texttt{sources/190819/source0.tex}.
\begin{theorem}[L2 Projection Approximation of SKM Algorithm]
\label{thm:appendix-skm-l2}
Consider the inequality constraint $Az\leq b$ where $A\in \mathbb{R}^{p\times n}$ and $b \in \mathbb{R}^p$, with feasible region $\mathcal{P}=\{z: Az\leq b\}$. Let the SKM method start from initial point $z_{0}$ and obtain $z_{k}$ after $k$ iterations. Then the expected distance between the iterate and initial point satisfies:
$$\mathbb{E}[d(z_{k},z_{0})] \leq 2 \cdot d(z_{0},\mathcal{P})$$
where $d(z_0, \mathcal{P}) = \min_{z \in \mathcal{P}} \|z - z_0\|_2$ is the distance from the initial point to the feasible region.
\end{theorem}

\begin{assumption}[Non-degeneracy]
\label{assump:nondeg}
There exists a constant $c > 0$ such that $\|a_i(x)\| \geq c$ for all $i,x$.
\end{assumption}

\begin{assumption}[Non-degeneracy of Tie Events]
\label{assump:tie}
For any fixed $\omega, k$ and any relevant constraint indices $i,j$, the functions $(a_i(x) - a_j(x))^T W_k(x,\omega) - (b_i(x) - b_j(x))$ and $a_i(x)^T W_k(x,\omega) - b_i(x)$ do not vanish identically on any open subset of input space $\mathbb{R}^{n_{\text{in}}}$.
\end{assumption}

\begin{theorem}[End-to-End Trainability of SKM Algorithm]
\label{thm:appendix-trainability}
Under Assumptions~\ref{assump:nondeg}--\ref{assump:tie}, the SKM algorithm has the following trainability properties:

\textbf{(i) Well-defined Expected Gradient}: $\mathbb{E}_\omega[W_K(x,\omega)]$ is almost everywhere differentiable with respect to $x$, and the expected gradient $\nabla_x \mathbb{E}_\omega[W_K(x,\omega)]$ is well-defined.

\textbf{(ii) Finite Gradient Variance}: $\text{Var}(\nabla_x W_K(x,\omega)) < \infty$.

\textbf{(iii) Unbiased Gradient Estimation}: The path derivative computed by automatic differentiation is an unbiased estimator of the expected gradient:
$$\mathbb{E}_\omega[\nabla_x W_K(x,\omega)] = \nabla_x \mathbb{E}_\omega[W_K(x,\omega)]$$
\end{theorem}

\begin{proof}
We establish the result through a series of lemmas.

\begin{lemma}[Boundedness of Iteration Sequence]
\label{lem:boundedness}
Based on the proof of Theorem~\ref{thm:appendix-skm-l2}, the SKM algorithm has Fejér monotonicity property. The iteration sequence $\{w_k\}$ satisfies: for any feasible point $w^* \in \mathcal{C} = \{w : A(x)w \leq b(x)\}$,
$$\|w_{k+1} - w^*\|^2 \leq \|w_k - w^*\|^2 - \delta(2-\delta)\frac{(r_{i^*}^+)^2}{\|a_{i^*}\|^2}$$
Therefore, there exists a constant $B$ such that $\sup_{k \geq 0} \|w_k\| \leq B < \infty$.
\end{lemma}

\begin{lemma}[Jacobian Properties of Single-step Mapping]
\label{lem:jacobian}
The Jacobian of single-step mapping $F(w,A,b,\mathcal{S})$ satisfies:

\textbf{(1) With respect to w}: When $a_{i^*}^T w > b_{i^*}$ and $\arg\max$ is unique:
$$\frac{\partial F}{\partial w} = I - \delta P_{a_{i^*}}, \quad \text{where} \quad P_a = \frac{aa^T}{\|a\|^2}$$
Since $P_a$ is a rank-1 projection matrix with eigenvalues $\{1, 0, \ldots, 0\}$:
$$\left\|\frac{\partial F}{\partial w}\right\|_2 = \max\{|1-\delta|, 1\} = 1 \quad (0 < \delta < 2)$$

\textbf{(2) With respect to b}: When constraints are active:
$$\left\|\frac{\partial F}{\partial b}\right\| \leq \frac{\delta}{c} := C_b$$

\textbf{(3) With respect to A}: Using Lemma~\ref{lem:boundedness} with $\|w\| \leq B$ and $|r_{i^*}^+| \leq MB + M$:
$$\left\|\frac{\partial F}{\partial A}\right\| \leq \delta \left( \frac{B}{c} + \frac{3(MB + M)}{c^2} \right) := C_A$$
\end{lemma}

\begin{lemma}[Zero Measure of Tie Events]
\label{lem:tie}
Under Assumption~\ref{assump:tie}, define the tie event
$$T_k = \{(x,\omega): \exists i \neq j \in \mathcal{S}_k, a_i^T W_k(x,\omega) - b_i = a_j^T W_k(x,\omega) - b_j\}$$
Then $\mu(T_k) = 0$, where $\mu$ is the Lebesgue measure.
\end{lemma}

\begin{proof}[Proof of Lemma~\ref{lem:tie}]
For fixed $\omega$, $W_k(x,\omega)$ is a piecewise differentiable function of $x$. The condition $a_i^T W_k - b_i = a_j^T W_k - b_j$ is equivalent to $(a_i - a_j)^T W_k = b_i - b_j$, which under Assumption~\ref{assump:tie} defines hypersurfaces in $x$ space. Similarly, the boundary where a single constraint transitions from non-violation to violation $a_i^T W_k = b_i$ also forms hypersurfaces.

Since $|\mathcal{S}_k| = \beta$ is finite, there are only finitely many constraint pairs $(i,j)$, thus $T_k$ is a union of finitely many hypersurfaces. Each hypersurface has zero Lebesgue measure in $n_{\text{in}}$-dimensional space, and a finite union of zero-measure sets remains zero-measure, hence $\mu(T_k) = 0$.

Similarly, $\mu(\bigcup_{k=0}^{K-1} T_k) = 0$ as it is a finite union of zero-measure sets.
\end{proof}

Now we prove the main theorem parts:

\textbf{(i) Well-defined Expected Gradient}: 
Through mathematical induction, at non-tie points we apply the chain rule:
$$\nabla_x W_{k+1} = \frac{\partial F}{\partial w}\nabla_x W_k + \frac{\partial F}{\partial A}\nabla_x A + \frac{\partial F}{\partial b}\nabla_x b$$

Let $C = C_A L_A + C_b L_b$. From Lemma~\ref{lem:jacobian}, we obtain the recurrence:
$$\|\nabla_x W_{k+1}\| \leq 1 \cdot \|\nabla_x W_k\| + C$$
Therefore: $\|\nabla_x W_k\| \leq kC$.

Setting $D_K = KC < \infty$, by the dominated convergence theorem and Lemma~\ref{lem:tie}:
$$\nabla_x \mathbb{E}_\omega[W_K(x,\omega)] = \mathbb{E}_\omega[\nabla_x W_K(x,\omega)]$$

\textbf{(ii) Finite Gradient Variance}: 
$$\mathbb{E}[\|\nabla_x W_K\|^2] \leq (KC)^2 < \infty$$

\textbf{(iii) Unbiased Gradient Estimation}: Since the probability $P(\omega)$ of sampling sequence $\omega$ is independent of parameter $x$:
$$\nabla_x \mathbb{E}_\omega[W_K(x,\omega)] = \sum_\omega P(\omega) \nabla_x W_K(x,\omega) = \mathbb{E}_\omega[\nabla_x W_K(x,\omega)]$$
\end{proof}

\end{document}
